\documentclass[12pt]{article}

% fonts

\usepackage[T1]{fontenc}
\usepackage{newtxtext}
\usepackage{bm}
\usepackage{amsthm} % before newtxmath b/c latex
\usepackage{newtxmath}
% \usepackage[bb=mt,cal=mt]{mathalpha}

\usepackage{amsmath}
\usepackage{amsfonts}
\usepackage{thmtools}
\usepackage{mathtools}

% spacing

\usepackage[parfill]{parskip}
\setlength{\parindent}{0pt}
\setlength{\parskip}{2ex plus 0.2ex minus 0.2ex}

% geometry of the page

\usepackage[
paper  = letterpaper,
left   = 1in,
right  = 1in,
top    = 1in,
bottom = 1in,
]{geometry}

% references

\usepackage{natbib}


\begin{document}

\begin{flushleft}
\textbf{Matrix factorization techniques for recommender systems (Koren et al., 2009)} \\
Satvik Saha (ss7228) \\
\today
\end{flushleft}

\vspace{0.1in}


This paper describes a recommender system, predicting \emph{users'} ratings for
\emph{items} (in this instance, films).
By assigning each user $u$ a latent factor $p_u \in \mathbb{R}^k$ and each item
$i$ a latent factor $q_i \in \mathbb{R}^k$, the rating $r_{ui}$ user $u$ has
for $i$ is modeled as approximately $q_i^\top p_u$.
If all user-item pairs $(u, i)$ were observed, or if we choose to impute
missing entries, recovering the latent factors could be presented as a matrix
factorization problem of the form $R = Q^\top P$.
Here, with only pairs from $\mathcal{K}$ observed, we may instead minimize a
squared error loss $\sum_{(u, i) \in \mathcal{K}} (r_{ui} - q_i^\top p_u)^2$.
This model can be extended by adding regularization terms, including user and
item biases, user and item attributes, and even temporal variations.

This has the usual topic modeling interpretation, with each item being a
$q_i$-mixture of $k$ topics, and each user having preferences for those topics
proportional to $p_u$.
The better aligned these vectors are, the higher the score $q_i^\top p_u$ for
the pair $(u, i)$.

Here, user attributes $Y$ and item attributes $X$ have been incorporated by
replacing the user latent vector $p_u$ with $p_u + (\text{item attributes for
items preferred by user $u$}) + (\text{user attributes})$.
A probabilistic model (such as a modified version of FLAIR) might use these
attributes in the following manner: \[
    r_{ui} \sim f(q_i^\top p_u + \lambda_u^\top x_i + \eta_i^\top y_u),
\] introducing more latent factors $\lambda_u, \eta_i$, or by introducing
hierarchy \begin{align*}
    r_{ui} \mid p_u, q_i &\sim f(\,\cdot\,\mid q_i^\top p_u), \\
    p_u &\sim \pi_1(\,\cdot\,; y_u), \\
    q_i &\sim \pi_2(\,\cdot\,; x_i).
\end{align*}
Probabilistic models like these would perhaps offer more information (say about
uncertainty) than point estimates obtained by direct matrix factorization.


\nocite{*}
\bibliographystyle{plainnat}

\vfill

\begingroup
\renewcommand{\section}[2]{}%
\textbf{References}
\bibliography{week03}
\endgroup


\end{document}

%%% Local Variables:
%%% mode: latex
%%% TeX-master: t
%%% End:
